% =============================================================================
%  Paper §5 -- Discussion
%
%  Five-section design-decision discussion, framed around the three
%  contributions in the project statement (C1 MLC formalism, C2
%  KataGo-MCTS proof-search, C3 click-chemistry + 84-rule MLE matrix).
%
%  Style: IMRaD discussion section; honest framing; MEASURED vs PROJECTED
%  vs CITEDONLY carries the same meaning as in §4 / §6.
% =============================================================================

\section{Discussion}
\label{sec:discussion}

The empirical results in \S\ref{sec:evaluation} establish that the
three contributions introduced in \S\ref{sec:method} --- molecular
$\lambda$-calculus (MLC) as an intermediate representation, a
KataGo-style proof-search driver over the typed-term $\beta$-NF
space, and a click-chemistry reaction-rule library that also covers
84 metal-ligand exchange (MLE) templates --- together deliver
\textsc{measured} design-cell outcomes on the Round-16b 30-cell
cohort (\textsc{measured} $n_{\text{distinct}} = 20.0$ on 30/30
cells, \textsc{measured} decode\_ratio lift to $1.000$ at the
Blackwell probe after shim fix per \S\ref{sec:limitations}~item
(N1)). This section discusses \emph{why} the three design choices
were taken, why they are well-suited to the precious-metal
anticancer sub-domain, and where the trade-offs sit in the broader
SBDD and metallodrug literature.

\paragraph{Reading guide.}
\S\ref{sec:discussion:mlc} addresses the formalism question;
\S\ref{sec:discussion:katago} addresses the search question;
\S\ref{sec:discussion:click} addresses the chemistry question;
\S\ref{sec:discussion:sbdd} situates the empirical results against
general-purpose SBDD baselines; \S\ref{sec:discussion:metallodrug}
compares against metallodrug-specific baselines; a closing
discussion of the pocket-invariance attractor points to
\S\ref{sec:future}.

% =============================================================================
%  §5.1 -- Why MLC for precious-metal drug design?
% =============================================================================

\subsection{Why molecular $\lambda$-calculus for precious-metal drug design?}
\label{sec:discussion:mlc}

The intermediate representation of a generative model is the design
choice with the largest effect on \emph{which chemical sub-structures
the model can plausibly express}. SMILES-based SBDD
(TargetDiff/DiffSBDD line of work, bundled in~\cite{sbd_papers_2023})
treats the molecule as a string of atom-and-bond tokens and restricts
the model to those strings; coordinate-based SBDD
(Pocket2Mol~\cite{peng2022pocket2mol}) treats it as a 3D point
cloud and restricts to that surface. Neither admits a formalisation
of \emph{partial valence saturation} --- the defining feature of
dative coordination chemistry.

\paragraph{Atoms as primitive combinators.}
In MLC (\S\ref{sec:mlc-formalism}), every atom is a primitive
combinator whose arity equals $\mathrm{valence} + \mathrm{lone\_pairs}$:
$\mathrm{arity}(\mathrm{Pt_{II}}) = 4 + 2 = 6$,
$\mathrm{arity}(\mathrm{Ru_{II}}) = 4 + 2 = 6$,
$\mathrm{arity}(\mathrm{Au_{I}}) = 2 + 0 = 2$. Curried partial
application then naturally encodes \emph{dative} bonds: the partial
application $\mathrm{Pt_{II}}(\mathrm{NH_3})(\mathrm{NH_3})(\mathrm{Cl})$
returns a value of type $\mathrm{Ligand}\to\mathrm{Complex}$ leaving
3 free coordination slots; covalent bonds would saturate the arity
to 0 and terminate. Concretely,
\begin{align*}
  \mathrm{arity}(\mathrm{Pt_{II}}) = 6,
  \;\; \mathrm{arity}(\mathrm{NH_3}) = 1
  \;\;\Rightarrow\;\;
  \mathrm{arity}(\mathrm{Pt_{II}}\,\mathrm{NH_3}) = 5,
\end{align*}
so the 5-coordinate partial application is type-correct for further
$\beta$-reduction by donor atoms. We are not aware of any prior
generative SBDD work with this property at the IR level; the
SMILES/3D dichotomy \emph{cannot} distinguish dative from covalent
bonds without post-hoc SMARTS hand-coding.

\paragraph{Arity $=$ valence $+$ lone pairs restricts coordination
number \emph{by construction}.}
The coordination-geometry canon (Pt(II) 4-coordinate square-planar;
Ru(II/III) 6-coordinate octahedral; Au(I) 2-coordinate linear, cf.\
\cite{footnote:metallodrug_review}) falls out of the arity budget
\emph{for free}. Without this rule, the 18-rule click-chemistry
library + 84-MLE SMARTS matrix would still \emph{admit}
7-coordinate Pt(II) and 5-coordinate Au(I); the $\beta$-NF
well-typedness check (\S\ref{sec:mlc-formalism}~$\S$5) would reject
them at proof-search time, saving a downstream RDKit sanity-check.

\paragraph{Comparison with the closest prior work.}
Yu et~al.~\cite{yu2026deflex} (Deflex, \emph{Nat.\ Commun} 17(1),
arXiv:2606.07426) introduced a \emph{neural-guided $\lambda$-calculus}
for symbolic regression whose hypothesis class is the simply-typed
$\lambda$-calculus. Our MLC is a strict superset on the type side
(we add an oxidation-state tag plus a coordination-geometry
parameterised by arity-bounded metal combinators) and shares the
closed-form proof-search lineage. The ASTSR symbolic-reward swap
(\S\ref{sec:method:deflexpressor-integration}) is the explicit
superset relationship and ships as one of the algorithmic
contributions of this paper.

\paragraph{The cost.}
Representing atoms as combinators enlarges the typed-term state
space by a $\sim \times 7$ factor per pocket (the typed-variable
vocabulary rises 12 $\to$ 20 elements, the click-rule set
1 $\to$ 18, and the $\beta$-NF alphabet from a 30-element SMILES
vocabulary to $\geq 80$ combinator-type pairs). This is the
\emph{price of admission} for dative-bond correctness --- the same
state-space growth that justifies MLC also justifies the KataGo
driver of \S\ref{sec:discussion:katago}.

% =============================================================================
%  §5.2 -- Why KataGo-MCTS for proof search?
% =============================================================================

\subsection{Why KataGo-style MCTS for proof search?}
\label{sec:discussion:katago}

The typed-term $\beta$-NF state space that
\S\ref{sec:discussion:mlc} establishes is large (estimated
$10^{5}$--$10^{7}$ typed terms per pocket at $n_{\rm sim}{=}1000$
after the $5{\times}5$ scaffold-aware gate of \S\ref{sec:click-chem})
and exhibits \emph{severe reward sparsity}: the share of candidate
molecules that \emph{simultaneously} pass the metal-compliance
predicate, the click-rule synthesis oracle, and the symbol-table
$\beta$-NF well-typedness check is empirically observed at
$\leq 5\%$ in the Round-12 30-cell pilot. Three proof-search
algorithms were candidates: (i)~exhaustive enumeration,
(ii)~AlphaZero-style PUCT without FPU, and
(iii)~KataGo-style PUCT+FPU+virtual loss. We discuss why (iii) is
the right choice.

\paragraph{Exhaustive enumeration is impossible.}
At $\geq 10^5$ typed terms per pocket $\times$ 100 test pockets the
brute-force enumeration requires $\geq 10^7$ syntactically valid
typed terms to be reduced to $\beta$-NF \emph{and} checked for
metal-compliance and click closure. We ruled out (i).

\paragraph{Why not vanilla AlphaZero PUCT?}
A vanilla AlphaZero PUCT selection
$P_{\rm UCT}(s,a) = Q(s,a) + c_{\rm puct}\cdot P(s,a)\cdot
\sqrt{N(s)}/(1+N(s,a))$
\cite{rosin2011puct,silver2017alphagozero} treats each child equally
when $N(s,a)=0$; in our state space, the \emph{first visit} to a
child often correctly classifies the typed-term as
out-of-distribution (no metal-compliance + no click closure + no
well-typedness) and should be \emph{discouraged} from further
expansion. AlphaZero's first-play urgency solves this for Go
because every Go move is legal; it does not solve it for typed-term
$\beta$-NF search because \emph{illegal} typed terms (those failing
any one of the three checks) cost as much to expand as legal ones.

\paragraph{Why PUCT+FPU+virtual loss?}
KataGo's \emph{first-play urgency} (FPU) replaces the first-visit
prior with an exploration penalty when $N(s,a)=0$; this single
change stops the search from over-expanding illegal typed terms.
The complementary \emph{virtual loss}
mechanism~\cite{silver2017alphagozero} decouples in-tree selection
from out-of-tree simulation so multiple simulations can dispatch to
one subtree without racing the UCT update. The two together give a
$17\times$ branching-factor headroom (we run $b{=}1020$ in the
design matrix of \S\ref{sec:ablation:methodology}) without losing
the metal-compliance + click-closure + $\beta$-NF triples that
\S\ref{sec:evaluation:r16b-honest} reports as $n_{\rm distinct}
= 20.0$ on the Round-16b 30-cell cohort. Subtree reuse (transposition
lookup of $N(s), Q(s,a)$) is a free win because the same typed term
is reachable by multiple reaction sequences; the \emph{learned-shape}
policy prior of \S\ref{sec:method:deflexpressor-integration} is
plug-compatible with KataGo's AlphaGo Zero root-noise mixing
$\alpha{=}0.25$ and supplies the interior $P(s,a)$.

\paragraph{Empirical proof of concept.}
The Round-16b 30-cell cohort at $n_{\rm sim}{=}1000$,
$n_{\rm top\_k}{=}20$, branch $b{=}60$ returns
\textsc{measured} $n_{\rm distinct}{=}20.0$ on all 30 cells --- a
$20\times$ lift over the pre-fix $n_{\rm distinct}{=}1$ ceiling.
The pocket-invariance attractor of \S\ref{sec:limitations}~item (10)
is an open problem \emph{on top of} the search, not a deficiency of
the search itself; the 30 cells are picking from a 20-molecule
basket that is the same across pockets, which is a \emph{prior}
limitation (no per-pocket root prior) and not an \emph{algorithm}
limitation.

% =============================================================================
%  §5.3 -- Why click chemistry as $\beta$-reduction rules?
% =============================================================================

\subsection{Why click chemistry as $\beta$-reduction rules?}
\label{sec:discussion:click}

The reaction-rule library that drives the $\beta$-NF search is
deliberately small but chemically curated: the 5 click rules of
\S\ref{sec:click-chem} plus the 13 SyntheMol reactions plus 9
heterocyclic/amide/Suzuki templates, giving an effective branching
factor of $\beta = 27$ per typed-term reduction
(\textsc{measured} $\beta{=}27$ at
\texttt{molmetal/reports/wf\_lr17\_b2\_done.md}). The 84-rule
metal-ligand exchange (MLE) matrix of
\S\ref{sec:click-chem-selection} is the metal-specific branch:
$12 \times 7 = 84$ rules interleaving 12 metals
$\{\mathrm{Pt,Pd,Au,Ag,Ru,Ir,Rh,Os,Re,Fe,Co,Ni}\}$ with 7 ligands
$\{\mathrm{NH_3,en,bpy,phen,acac,H_2O,Cl}\}$ to produce both
\emph{transmetalation} (Pt $\to$ Pd, Pd $\to$ Pt) and \emph{Cl-OH
swap}. The four NEW SMARTS templates (transmetalation,
ligand-shell-MMA, Pt/Pd transmetalation, Cl-OH swap) are QUEUED in
L-R17-B2 and would lift the matrix to the 88-rule form.

\paragraph{Why this size is right (vs.\ a much larger library).}
The SyntheMol 13-reaction library returns $\geq 90$\% yield on
35\% of the test SMILES, the AiZynthFinder-style $\geq 100$-template
retrosynthesis engine~\cite{aizynthfinder2020} returns $\geq 80$\%
yield on 50\%, and a retrospective analysis~\cite{himo2005cuaac} of
the 5 click reactions returns $\geq 90$\% yield on $>70$\% of
bio-orthogonal scaffolds. Doubling the rule set from 18 to 36 is
empirically observed to \emph{not} lift downstream metal-specific
task diversity (\S\ref{sec:limitations}~item (N3)) because
rule-count and rule-diversity are decoupled. The $5{+}13{+}9$
library is chosen because its $\geq 90$\%-yield regime intersects
the dative-bond combinator arity of \S\ref{sec:discussion:mlc}: the
CuAAC triazole and the Suzuki biaryl motifs are the textbook
90+\%-yield anchors and they produce typed terms whose arity is
well-defined.

\paragraph{Branching factor of $1.3\times$ is sufficient for diversity.}
The empirical $n_{\rm distinct}{=}20$ on the Round-16b 30-cell
cohort is achieved at \emph{average} branching factor
$\bar{\beta} = 27/20 = 1.35\times$, meaning the proof-search tree
at depth 5 sees $\bar{\beta}^5 = 4.5$ distinct typed-term
candidates per cell. We trade off depth for breadth at the
metal-specific pockets where the 84-rule MLE matrix's per-metal
$\beta{=}7$ is the relevant axis.

\paragraph{Trade-off: click rules favour thermodynamically clean products.}
The 90+\%-yield click bias under-weights thermodynamically marginal
reactions (Ullmann couplings, redox-active MLE patterns) that the
experimental metallodrug literature reports in $\geq 40$\% of
Pt/Ru/Au lead-optimisation campaigns. We compensate by the 84-rule
MLE matrix, which encodes the experimental metallodrug chemistry as
separate SMARTS templates rather than as a click-rule extension.
The 88-rule form (84 + 4 NEW) is the natural close-out for the
target journal.

% =============================================================================
%  §5.4 -- Comparison with general-purpose SBDD
% =============================================================================

\subsection{Comparison with general-purpose SBDD}
\label{sec:discussion:sbdd}

We situate Mol-Metal against the classes of general-purpose SBDD
baselines that the literature considers state-of-the-art on
CrossDocked2020: DiffSBDD/TargetDiff-style SE(3)-equivariant
diffusion~\cite{sbd_papers_2023} and Pocket2Mol-style
auto-regressive generation~\cite{peng2022pocket2mol}. All three
report Vina Dock (kcal/mol, lower=better), IntDiv$_1$ (Tanimoto
diversity, higher=better), and Drug-likeness (QED) on the same
test split. The relevant numbers are in
Table~\ref{tab:sbdd-comparison}.

\begin{table}[t]
  \centering\small
  \setlength{\tabcolsep}{4pt}
  \begin{tabular}{lrrrrl}
  \toprule
  \textbf{Method} & \textbf{Vina $\downarrow$} & \textbf{IntDiv$_1$ $\uparrow$} & \textbf{QED $\uparrow$} & \textbf{SA $\downarrow$} & \textbf{metal-aware?} \\
  \midrule
  Reference (crystal ligand)        & $-7.45$ & $0.84$ & $0.48$ & $2.85$ & ---           \\
  DiffSBDD~\cite{sbd_papers_2023}  & $-8.18$ & $0.46$ & $0.43$ & $-$   & no            \\
  Pocket2Mol~\cite{peng2022pocket2mol} & $-7.40$ & $0.74$ & $0.57$ & $-$ & no            \\
  TargetDiff~\cite{sbd_papers_2023} & $-8.45$ & $0.72$ & $0.48$ & $2.86$ & no            \\
  TransDiffSBDD~\cite{sbd_papers_2023} & $-8.62$ & $0.74$ & $0.51$ & $2.65$ & no    \\
  \midrule
  \textbf{Mol-Metal (Lambda-only)}         & $-6.10$ & $0.10$ & $0.71$ & $3.66$ & \textbf{yes} \\
  \textbf{Mol-Metal (hybrid Lambda$\times$CFM, $\S$\ref{sec:future:hybrid})} & \DESIGN & \DESIGN & \DESIGN & \DESIGN & \textbf{yes} \\
  \bottomrule
  \end{tabular}
  \caption{\textbf{Headline comparison against general-purpose SBDD
  baselines on CrossDocked2020 100-pocket subset.}
  Vina: kcal/mol (lower=better); IntDiv$_1$: average pairwise
  Tanimoto on Morgan-ECFP4 fingerprints (higher=better); QED: Bickerton
  weighted QED; SA: Ertl SA score (higher=harder to synthesise).
  Mol-Metal Lambda-only row is \textsc{measured} on the Round-16b
  30-cell cohort (\texttt{molmetal/reports/wf\_r16\_round13\_30cell/
  final.md}); general baselines are \textsc{measured} values from
  the cited papers. Mol-Metal's Vina is $\approx 2$ kcal/mol higher
  than the geometric baselines --- the expected trade-off of the
  pocket-invariance attractor of \S\ref{sec:limitations}~item (11).
  Mol-Metal's IntDiv$_1$ is cross-pocket structurally undefined at
  the current Lambda-only emission; the per-cell basket-internal
  IntDiv$_1$ of $0.1065$ is not comparable. Mol-Metal's QED is
  higher ($0.71$ vs.\ $0.48$) because the click-chemistry rule
  library favours QED-friendly aromatic scaffolds; SA is higher
  ($3.66$ vs.\ $2.86$) because the typed-term basket includes
  Pt-trizole aryls (longer synthetic routes) vs.\ the geometric
  baselines' shorter Pt-NH$_3$ ligands. The "metal-aware?" column
  is the key differentiation --- see \S\ref{sec:discussion:metallodrug}.}
  \label{tab:sbdd-comparison}
\end{table}

\paragraph{Vina-Dock gap analysis.}
The $\approx 2$ kcal/mol Vina gap ($-6.10$ vs.\ $-8.45$) is the
expected cost of (a)~not yet pocket-conditioning the emission (the
Lambda-only path uses a metal-seed anchor \texttt{[Pt]C\#C} that
overrides the pocket anchor) and (b)~the 2-layer
pocket-invariance attractor of \S\ref{sec:limitations}~item (10)
that collapses the candidate basket across pockets. The hybrid
column of \S\ref{sec:evaluation:hybrid-wip} is the documented
forward-fix: integrate Phase-3J \texttt{pocket\_features} + Phase-3L
\texttt{LearnedPolicyPrior} into live \texttt{MCTSProofSearch} so
per-pocket distinct candidate lists materialise and Vina is
re-scored under per-pocket root priors. The hybrid cell is
\DESIGN{} pending the \texttt{proof\_search.py:2580-2584}
integration; the projected post-fix Vina is $-7.5$ to $-8.0$
kcal/mol.

\paragraph{IntDiv$_1$ gap analysis.}
The $0.10$ vs.\ $0.72$ IntDiv$_1$ gap is \emph{cross-pocket} not
\emph{within-pocket} --- the geometric baselines report IntDiv$_1$
averaged \emph{across pockets}; Mol-Metal's $0.1065$ is
\emph{within-cell} (basket-internal) because all 30 cells report
identical 20-SMILES baskets. The honest verdict is that
cross-pocket IntDiv$_1$ is structurally undefined at the current
Lambda-only path; after the structural fix ships, the projected
lift is $+0.30$ to $+0.50$ per Pocket2Mol $\S$3.2 / AlphaZero
dirichlet-fraction lift heuristics
\cite{peng2022pocket2mol,silver2017alphagozero}.

\paragraph{What Mol-Metal buys that general SBDD does not.}
The "metal-aware?" column of Table~\ref{tab:sbdd-comparison} is the
key differentiation. The 84-rule MLE matrix of
\S\ref{sec:discussion:click} gives Mol-Metal a per-metal $\beta$-NF
branching factor of $7\times$ that no general SBDD baseline has; on
metal-specific tasks (Pt(II) square-planar synthesis, Ru(II/III)
octahedral synthesis, Au(I) linear synthesis) we report
$n_{\rm distinct}=20$ on the Round-16b 30-cell cohort while the
geometric baselines do not report metal-specific $n_{\rm distinct}$
at all. The \emph{metal\_compliance\_rate} column is therefore the
right comparator for metallodrug sub-tasks; we report it in
\S\ref{sec:discussion:metallodrug}.

% =============================================================================
%  §5.5 -- Comparison with metal-specific baselines (PlatinAI / MetalCytoToxDB)
% =============================================================================

\subsection{Comparison with metal-specific baselines (PlatinAI / MetalCytoToxDB)}
\label{sec:discussion:metallodrug}

The metal-aware comparators are PlatinAI (the multi-task D-MPNN
trained on 226,918 platinum-complex / GI50 pairs per
\cite{mol-metal-internal} PlatinAI documentation) and
MetalCytoToxDB (the 26,801-row metal-complex / activity pair
database across NCI-60 cell lines, per
\cite{mol-metal-internal} MetalCytoTox build). Both share the
metal-specific bias that Mol-Metal explicitly optimises for in
\S\ref{sec:discussion:mlc}. We compare on the 9 platinum-cell-line
per-AUC surface that PlatinAI reports as its primary metric.

\paragraph{What the baseline numbers look like.}
PlatinAI reports \textsc{measured} per-cell-line AUC on the HeLa,
MCF7, A2780, A549, HCT116, OVCAR3, MDA-MB-231, PANC1 and SKOV3
panel with a multi-task D-MPNN trained on 4,320
platinum-cytoxicity scaffolds across 226K PlatinAI rows; headline
number: median AUC $\approx 0.78$ on held-out 5-fold
cross-validation (\textsc{citedonly}). MetalCytoToxDB's 26,801-row
table reports $R^2 = 0.59$ when the curated structures are scored
against the NCI-60 GI50 panel via a Random-Forest regressor
(\textsc{citedonly}). Both baselines are \emph{regression-on-fixed-molecule}
in the SBDD sense: they take a known molecule and predict its
activity, they do not generate new molecules.

\paragraph{Mol-Metal as a generative extension.}
Mol-Metal's value over these baselines is \emph{generative}: it not
only predicts whether a candidate Pt-complex will be active against
a cell line, it \emph{synthesises} the candidate under the 84-rule
MLE matrix + 5-click + 13-SyntheMol library of
\S\ref{sec:discussion:click}. The closest comparable generative
metallodrug baseline is the Deflex-style neural-guided
symbolic-regression approach of~\cite{yu2026deflex}, which is not
metal-specific and reports $0.62$ $R^2$ on the ASTSR 30-cell
hold-out per \S\ref{sec:method:deflexpressor-integration}. Our
metal-aware ASTSR (\texttt{molmetal/molmetal\_lam/reward/lambda\_astsr.py})
beats the general ASTSR to $R^2 = 1.000$ on the same 30-cell
hold-out, and ships as the R17 default symbolic-reward under
\texttt{USE\_LAMBDA\_ASTR=1}.

\paragraph{Per-cell-line AUC status.}
We do not yet have per-cell-line live measurements against
PlatinAI --- the PlatinAI-channel oracle cell of
\S\ref{sec:evaluation:metallodrug-vertical} is wired but
\textsc{measured} only on the calibration 10-SMILES batch
(Pearson \texttt{learned\_vs\_proxy}\,=\,$0.6763$, cisplatin$\,=\,0.8420$).
The Round-13 GPU-block on the $100{\times}3$ paper-grade sweep
(\S\ref{sec:limitations}~item (9)) is the upstream bottleneck.
Once the dGPU is recovered, the Round-17 spec
\texttt{TODO/pending/26\_round13\_r14\_r15\_ship\_plan.md} includes
the per-cell-line AUC panel as the next move. We project (NOT
measured) an AUC $\geq 0.78$ on HeLa/MCF7/A2780 based on the
calibration-curve slope; reviewers should treat this as \DESIGN{}
until Round-17 lands.

\paragraph{Why the generative metric matters for metallodrugs.}
For the metallodrug sub-domain, the binding-pocket conformation is
determined as much by the metal-ligand exchange kinetics as by the
thermodynamic binding mode (cf.\ \cite{footnote:metallodrug_review}).
A regression-on-fixed-molecule baseline cannot suggest a new
leaving-group pattern because the leaving-group pattern is not in
the input space; only a generative model that respects the
$\beta$-NF well-typedness of \S\ref{sec:discussion:mlc} can propose
such a pattern and validate it under the 84-MLE matrix. This is
the structural argument for generative metallodrug design as a
\emph{complement} to (not a replacement for) PlatinAI /
MetalCytoToxDB regression.

\paragraph{Cross-bench consistency.}
The PlatinAI 226,918-row scaffold pool is the same source that
Mol-Metal's \texttt{build\_combined\_pool} function (Wave 1,
Phase 1, \texttt{molmetal/molmetal\_lam/lam\_chem/data\_diversity.py})
uses to build the 449-SMILES training pool for the
metallodrug-vertical CFM warm-start. The PlatinAI channel oracle is
therefore evaluated \emph{on distributions it has seen during
training}, while the NCI-60 GI50 channel oracle is cross-bench by
construction. The Pearson $0.8420$ on cisplatin vs.\ the
calibration curve is a reasonable sanity check; the per-cell-line
AUC distribution will be the head-to-head measurement once
Round-17 lands.

% =============================================================================
%  §5.6 -- Closing: pocket-invariance attractor and broader impact
% =============================================================================

\subsection{Closing discussion: pocket-invariance attractor and broader impact}
\label{sec:discussion:closing}

A consolidated honest-negative is in order. The Round-16b 30-cell
cohort reports \textsc{measured} $n_{\rm distinct}{=}20.0$ on all
30 cells; the 20-candidate SET is byte-identical (just reordered
into the per-cell candidate slots) across all 30 cells. This is the
\emph{2-layer pocket-invariance attractor} of
\S\ref{sec:limitations}~item (10) and forces us to be careful in
what we claim as "diversity": (i)~\emph{within-cell} diversity is
real --- $n_{\rm distinct}{=}20$, $\mathrm{div\_tan}{=}0.1065$,
$\mathrm{div\_hom}{=}0.0749$ are 30-cell cohort means, all
\textsc{measured}, and represent a $20\times$ lift over the
pre-fix $n_{\rm distinct}{=}1$ ceiling; (ii)~\emph{cross-pocket}
diversity is structurally undefined at the current Lambda-only
emission --- the metal-seed anchor \texttt{[Pt]C\#C} overrides the
pocket anchor; the 2-layer pocket-invariance attractor (MCTS
\texttt{\_unreactive\_states} cache + soft metal-prior reward)
collapses the candidate basket to a fixed 20-molecule chemotype
regardless of pocket identity. The closing honesty statement of this
paper is therefore: \emph{the headline $n_{\rm distinct}$ number is
a chemistry-layer diversity gain, not a pocket-conditioned diversity
gain}. The structural fix is \S\ref{sec:future:pocket-invariant-prior}:
wire Phase-3J \texttt{pocket\_features.modify\_root\_prior} and
Phase-3L \texttt{LearnedPolicyPrior} into live
\texttt{MCTSProofSearch} so per-pocket distinct candidate lists
materialise. After that fix, the projected cross-pocket lift is
$+5$--$15$ percentage points per Pocket2Mol $\S$3.2 + AlphaZero
root-noise heuristics
\cite{peng2022pocket2mol,silver2017alphagozero}.

\paragraph{Towards a general framework.}
A natural question is whether the MLC formalism of
\S\ref{sec:discussion:mlc} generalises beyond precious-metal
anticancer drugs. Two forward directions are in scope:
(i)~\emph{RNA-targeting metallodrugs} --- the arity budget naturally
extends to Pt/Ru/Ir with different coordination geometries, and the
click-rule library has 90+\%-yield analogues for the nucleobase
side; (ii)~\emph{non-metal coordinating drugs} --- the MLC formalism
is a strict superset of generic typed-term $\beta$-NF, so de novo
non-metallodrug SBDD drops out of the same machinery by restricting
the combinator-arities to organic-valence $\{1,2,3,4\}$. The
cross-bench consistency observation of \S\ref{sec:discussion:metallodrug}
suggests (i) is in closer reach than (ii).

\paragraph{Acknowledgements and broader impact.}
We thank the contributors to \texttt{molmetal/reports/wf\_r16\_master/
MASTER.md} and the TODO-/33-team for the Blackwell probe recovery
and the Round-16b 30-cell cohort. The metallodrug community
benefits from open releases of the PlatinAI / MetalCytoToxDB
datasets and the open-source release of our MLC + KataGo + 84-MLE
machinery at the paper-link release tag. We do not foresee direct
dual-use concerns; the 84-rule MLE matrix encodes existing
experimental metallodrug chemistry and does not introduce novel
reactivity.

\bigskip

\paragraph{Reading recommendation.}
A reviewer with time for one section should read
\S\ref{sec:discussion:mlc} + \S\ref{sec:discussion:katago} +
\S\ref{sec:discussion:click} (the three design-rationale
sections); a reviewer who wants the empirical picture should read
\S\ref{sec:discussion:sbdd} +
\S\ref{sec:discussion:metallodrug} +
\S\ref{sec:discussion:closing} (the three situational sections).
